\unpart{The Relativistic Reduction and Common-Invariant Program}
\section{From Event Composition to Causal Order}\label{sec:causalorder}
An occurrence graph has completed events as vertices and dependencies as arrows: a propagated consequence or retained record of one occurrence is an input to another. If this graph is acyclic, its transitive closure defines causal precedence. Recurrent carrier configurations can appear at distinct occurrences without making the occurrence graph cyclic.

Independent occurrences can be incomparable. A global enumeration is therefore additional bookkeeping unless an operational synchronization is supplied. Unit adjacent-step locality gives a finite support cone in graph distance versus event depth; a Lorentzian interpretation additionally requires isotropy, scaling, and transformations between realized clocks \citep{arrighi2019,sorkin2003}.

\section{Relational Distance and Emergent Metric Structure}\label{sec:metric}
A metric reduction must specify intervals, dimensional neighborhood structure, a volume measure, and operational clock and rod comparisons. Field distance and registered winding supply distinct inputs. Their combination becomes an interval only after compatible domains, units, and transformation laws are established.

For a sequence of fields $\mathcal F_\varepsilon$, the reduction should identify coarse observations and prove their convergence independently of microscopic relabeling. Proper time requires a future-directed physical clock law; reversal-odd cycle winding alone does not supply its elapsed duration or calibration.

\section{Retained History and Effective Geometry}\label{sec:historygeometry}
\setcounter{theorem}{0}
\begin{hypothesis}[History-conditioned accessibility]
In an appropriate transducer field, some retained relational distinctions affect subsequent admissible handoffs. After coarse-graining, this dependence admits a geometric description of causal accessibility.
\end{hypothesis}

The first obligation is an explicit residue-dependent admission or response law; the second is its geometric scaling. A proposed geometric source must be invariant under redundant record copying and consistent under reversible redistribution of information. Record length alone is not an admissible source.

Finite transport maps already support holonomy. A spacetime connection requires a frame or tangent interpretation, an operational loop scale, and a controlled relation between small-loop holonomy and curvature.

\section{General Relativity as a Coarse-Grained Transducer Description}\label{sec:gr}
The proposed geometric reduction retains causal organization while coarse-graining internal registration. Its target includes an operational metric, matter observables, and effective dynamics. In conventional units one target is
% Original equation number(s) 46; expression unchanged.
\setcounter{equation}{45}
\begin{equation}
G_{\mu\nu}+\Lambda g_{\mu\nu}
=\frac{8\pi G_N}{c_{\mathrm{SI}}^4}\,T_{\mu\nu}.
\end{equation}

A successful correspondence must specify the stress-energy map, coupling, approximation regime, corrections, and compatibility with covariant conservation \citep{carroll1997}. Local inertial behavior and Lorentz symmetry require their own argument; a discrete carrier alone supplies neither a preferred-frame violation nor an acceptable Lorentzian limit \citep{bombelli2009}. The present finite results fix the architectural inputs and consistency tests, not a continuum field equation.
